\documentclass[10pt,a4paper]{amsart}
\usepackage{amsmath,amssymb,amsthm,graphicx}
% Retained minimal preamble for the Stacks Project chapter more-algebra.tex (More on Algebra).
% The chapter file pinned-original-source.tex carries only \input{preamble} on its first line, so no
% preamble text lives inside the chapter itself.  The Stacks Project uses one single shared preamble for
% every chapter; the snapshot stages it identically as sources/stacks-chow/preamble.tex and
% sources/stacks-etale-cohomology/preamble.tex.  The block below is copied mechanically, line for line,
% from sources/stacks-chow/preamble.tex lines 14-18 (the Xy-pic setup) and 154-217 (the native theorem
% declarations together with the mathematical macros).  Nothing was added, removed or redefined; the
% wrapper class and the amsmath/amssymb/amsthm/graphicx base packages are supplied by the assembler.
% Chapter: more-algebra.tex -- More on Algebra.
% For commutative diagrams we use Xy-pic
\usepackage[all]{xy}

% We use 2cell for 2-commutative diagrams.
\xyoption{2cell}

% Theorem environments.
%
\theoremstyle{plain}
\newtheorem{theorem}[subsection]{Theorem}
\newtheorem{proposition}[subsection]{Proposition}
\newtheorem{lemma}[subsection]{Lemma}

\theoremstyle{definition}
\newtheorem{definition}[subsection]{Definition}
\newtheorem{example}[subsection]{Example}
\newtheorem{exercise}[subsection]{Exercise}
\newtheorem{situation}[subsection]{Situation}

\theoremstyle{remark}
\newtheorem{remark}[subsection]{Remark}
\newtheorem{remarks}[subsection]{Remarks}

\numberwithin{equation}{subsection}

% Macros
%
\def\lim{\mathop{\mathrm{lim}}\nolimits}
\def\colim{\mathop{\mathrm{colim}}\nolimits}
\def\Spec{\mathop{\mathrm{Spec}}}
\def\Hom{\mathop{\mathrm{Hom}}\nolimits}
\def\Ext{\mathop{\mathrm{Ext}}\nolimits}
\def\SheafHom{\mathop{\mathcal{H}\!\mathit{om}}\nolimits}
\def\SheafExt{\mathop{\mathcal{E}\!\mathit{xt}}\nolimits}
\def\Sch{\mathit{Sch}}
\def\Mor{\mathop{\mathrm{Mor}}\nolimits}
\def\Ob{\mathop{\mathrm{Ob}}\nolimits}
\def\Sh{\mathop{\mathit{Sh}}\nolimits}
\def\NL{\mathop{N\!L}\nolimits}
\def\CH{\mathop{\mathrm{CH}}\nolimits}
\def\proetale{{pro\text{-}\acute{e}tale}}
\def\etale{{\acute{e}tale}}
\def\QCoh{\mathit{QCoh}}
\def\Ker{\mathop{\mathrm{Ker}}}
\def\Im{\mathop{\mathrm{Im}}}
\def\Coker{\mathop{\mathrm{Coker}}}
\def\Coim{\mathop{\mathrm{Coim}}}

% Boxtimes
%
\DeclareMathSymbol{\boxtimes}{\mathbin}{AMSa}{"02}

%
% Macros for moduli stacks/spaces
%
\def\QCohstack{\mathcal{QC}\!\mathit{oh}}
\def\Cohstack{\mathcal{C}\!\mathit{oh}}
\def\Spacesstack{\mathcal{S}\!\mathit{paces}}
\def\Quotfunctor{\mathrm{Quot}}
\def\Hilbfunctor{\mathrm{Hilb}}
\def\Curvesstack{\mathcal{C}\!\mathit{urves}}
\def\Polarizedstack{\mathcal{P}\!\mathit{olarized}}
\def\Complexesstack{\mathcal{C}\!\mathit{omplexes}}
% \Pic is the operator that assigns to X its picard group, usage \Pic(X)
% \Picardstack_{X/B} denotes the Picard stack of X over B
% \Picardfunctor_{X/B} denotes the Picard functor of X over B
\def\Pic{\mathop{\mathrm{Pic}}\nolimits}
\def\Picardstack{\mathcal{P}\!\mathit{ic}}
\def\Picardfunctor{\mathrm{Pic}}
\def\Deformationcategory{\mathcal{D}\!\mathit{ef}}
\usepackage{cleveref}
\crefname{originallemma}{original lemma}{original lemmas}
\Crefname{originallemma}{Original lemma}{Original lemmas}
\makeatletter
\newlabel{algebra-lemma-cokernel-flat}{{00O0}{source}{}{sourcepointer-1}{}}
\newlabel{algebra-lemma-cokernel-flat@cref}{{[originallemma][0][]00O0}{[1][0][]source}{}{}{}}
\makeatother
\crefname{originallemma}{original lemma}{original lemmas}
\Crefname{originallemma}{Original lemma}{Original lemmas}
\makeatletter
\newlabel{lemma-splitting-unique}{{0BCG}{source}{}{sourcepointer-2}{}}
\newlabel{lemma-splitting-unique@cref}{{[originallemma][0][]0BCG}{[1][0][]source}{}{}{}}
\makeatother
\crefname{originalremark}{original remark}{original remarks}
\Crefname{originalremark}{Original remark}{Original remarks}
\makeatletter
\newlabel{derived-remark-truncation-distinguished-triangle}{{08J5}{source}{}{sourcepointer-3}{}}
\newlabel{derived-remark-truncation-distinguished-triangle@cref}{{[originalremark][0][]08J5}{[1][0][]source}{}{}{}}
\makeatother
\title{Cutting a pseudo-coherent complex in two pieces after localizing, once a cohomology group is generated by its reduction at a prime}
\author{The Stacks Project Authors}
\date{stacks-project:more-algebra.tex:local-snapshot-20261012; GNU FDL 1.2 or later}
\begin{document}
\maketitle
\setcounter{subsection}{1}
\noindent\textit{Editorial note.} All delivered mathematical text is the author's own, copied from the pinned Stacks Project chapter file more-algebra.tex; no statement or proof step was rewritten, reordered or removed. The counted statement is the labeled lemma environment at source lines 19743-19763 (canonical tag 0A1U) and the counted proof is the proof environment at source lines 19765-19819, which immediately follows it and is its own proof.\par\medskip
\noindent\textit{Editorial note.} The proof refers once to a statement of this chapter and twice to results of other chapters, namely the localization statement for cokernels and the distinguished triangle of truncations; all three are carried as navigation stubs with their roles recorded in the reasons.\par\medskip
\noindent\textit{Editorial note.} The statement and proof use the truncation notation tau of the shared preamble, which the retained preamble supplies verbatim.\par\medskip
\noindent\textit{Editorial note.} The chapter file pinned-original-source.tex carries only an input of the shared preamble on its first line, so no preamble text lives inside the chapter. The retained file sources/stacks-more-algebra/retained-preamble.tex is a mechanical, line-for-line copy of the project's single shared preamble as staged in the snapshot (Xy-pic setup plus the native theorem declarations and mathematical macros). Nothing was added, removed or redefined, and the excerpt serves only as preamble, never as a span coordinate.\par\medskip
\noindent\textit{Editorial note.} The source URL is the source-level Stacks Project URL, as required when more than one card is produced from one source; the per-claim tag is carried by claim\_key and by this note. The counted statement and proof cite no bibliography entry, so this card needs no bibliography.\par\medskip
\noindent\textit{Allowed standard prerequisites (editorial).}\begin{itemize}
\item A pseudo-coherent complex is represented by a bounded above complex of finite free modules, and a finite module over a local ring that is generated by its reduction can be generated by lifts modulo an arbitrarily high power of the maximal ideal. used at proof lines 19767-19768 and 19786-19794. Justifies the choice of the elements y\_b and the localization that makes the three listed conditions hold.
\item If a map of finite free modules has a surjective reduction modulo a prime and its cokernel becomes finite projective after inverting some element outside the prime, then the cokernel itself is finite projective after that localization. invoked at proof lines 19793-19799. Produces the finite projective quotient E\textasciicircum{}\{i+1\} that makes the upper truncation perfect.
\item For a distinguished triangle whose first term has no cohomology above a given degree and whose third term is perfect with amplitude starting at that degree, the triangle splits and exhibits the middle term as a direct sum of the two truncations. invoked at proof lines 19810-19818. Gives the displayed canonical isomorphism and finishes the proof.
\end{itemize}\medskip
\ifdefined\hypertarget\hypertarget{sourcepointer-1}{}\fi
\noindent\textit{Source reference.} Original lemma 00O0: sources/stacks-chow/tags.txt, source lines 878--878. The reference is the localization statement applied at proof lines 19793-19794 to obtain the three conditions after inverting an element outside the prime; the three conditions and their use are inside the counted span.\par\medskip
\ifdefined\hypertarget\hypertarget{sourcepointer-2}{}\fi
\noindent\textit{Source reference.} Original lemma 0BCG: sources/stacks-more-algebra/pinned-original-source.tex, source lines 19704--19704. The reference is applied at proof lines 19811-19812 to deduce the canonical isomorphism from the distinguished triangle; the identification of the upper truncation as perfect is inside the counted span.\par\medskip
\ifdefined\hypertarget\hypertarget{sourcepointer-3}{}\fi
\noindent\textit{Source reference.} Original remark 08J5: sources/stacks-chow/tags.txt, source lines 10838--10838. The reference is the standard distinguished triangle relating a complex to its two truncations, cited at proof lines 19817-19818; it is a navigation pointer to the triangle and the splitting statement that consumes it is inside the counted span.\par\medskip
\begin{lemma}
\label{lemma-better-cut-complex-in-two}
Let $R$ be a ring. Let $\mathfrak p \subset R$ be a prime ideal.
Let $K^\bullet$ be a pseudo-coherent complex of $R$-modules.
Assume that for some $i \in \mathbf{Z}$ the map
$$
H^i(K^\bullet) \otimes_R \kappa(\mathfrak p)
\longrightarrow
H^i(K^\bullet \otimes_R^{\mathbf{L}} \kappa(\mathfrak p))
$$
is surjective. Then there exists an $f \in R$, $f \not \in \mathfrak p$
such that $\tau_{\geq i + 1}(K^\bullet \otimes_R R_f)$ is a perfect
object of $D(R_f)$ with tor amplitude in $[i + 1, \infty]$ and
a canonical isomorphism
$$
K^\bullet \otimes_R R_f \cong
\tau_{\leq i}(K^\bullet \otimes_R R_f) \oplus
\tau_{\geq i + 1}(K^\bullet \otimes_R R_f)
$$
in $D(R_f)$.
\end{lemma}

\begin{proof}
In this proof all tensor products are over $R$ and we write
$\kappa = \kappa(\mathfrak p)$. We may assume that $K^\bullet$
is a bounded above complex of finite free $R$-modules. Let us
inspect what is happening in degree $i$:
$$
\ldots \to K^{i - 1} \xrightarrow{d^{i - 1}} K^i \xrightarrow{d^i}
K^{i + 1} \to \ldots
$$
Let $0 \subset V \subset W \subset K^i \otimes \kappa$ be defined
by the formulas
$$
V = \Im\left(
K^{i - 1} \otimes \kappa \to K^i \otimes \kappa
\right)
\quad\text{and}\quad
W = \Ker\left(
K^i \otimes \kappa \to K^{i + 1} \otimes \kappa
\right)
$$
Set $\dim(V) = r$, $\dim(W/V) = s$, and $\dim(K^i \otimes \kappa/W) = t$.
We can pick $x_1, \ldots, x_r \in K^{i - 1}$ which map by $d^{i - 1}$
to a basis of $V$. By our assumption we can pick
$y_1, \ldots, y_s \in \Ker(d^i)$ mapping to a basis of $W/V$.
Finally, choose $z_1, \ldots, z_t \in K^i$ mapping to a basis of
$K^i \otimes \kappa/W$. Then we see that the elements
$d^i(z_1), \ldots, d^i(z_t) \in K^{i + 1}$ are linearly independent
in $K^{i + 1} \otimes_R \kappa$.
By Algebra, Lemma \ref{algebra-lemma-cokernel-flat} we may after replacing
$R$ by $R_f$ for some $f \in R$, $f \not \in \mathfrak p$ assume that
\begin{enumerate}
\item $d^i(x_a), y_b, z_c$ is an $R$-basis of $K^i$,
\item $d^i(z_1), \ldots, d^i(z_t)$ are $R$-linearly independent in
$K^{i + 1}$, and
\item the quotient $E^{i + 1} = K^{i + 1}/\sum Rd^i(z_c)$ is finite projective.
\end{enumerate}
Since $d^i$ annihilates $d^{i - 1}(x_a)$ and $y_b$, we deduce from
condition (2) that $E^{i + 1} = \Coker(d^i : K^i \to K^{i + 1})$.
Thus we see that
$$
\tau_{\geq i + 1}K^\bullet =
(\ldots \to 0 \to E^{i + 1} \to K^{i + 2} \to \ldots)
$$
is a bounded complex of finite projective modules sitting in degrees
$[i + 1, b]$ for some $b$. Thus $\tau_{\geq i + 1}K^\bullet$ is perfect
of amplitude $[i + 1, b]$. Since $\tau_{\leq i}K^\bullet$ has no
cohomology in degrees $> i$, we may apply Lemma \ref{lemma-splitting-unique}
to the distinguished triangle
$$
\tau_{\leq i}K^\bullet \to K^\bullet \to \tau_{\geq i + 1}K^\bullet \to
(\tau_{\leq i}K^\bullet)[1]
$$
(Derived Categories, Remark
\ref{derived-remark-truncation-distinguished-triangle}) to conclude.
\end{proof}

\end{document}
